\scanpage{12}
prime to the residue characteristic, then the ring $\deleted{B=}$
\[
B=A\ T_1,..,T_r\ /\ (T_1^{n_1}-x_1,\ldots,T_r^{n_r}-x_r)
\]
% No brackets are visible around T_1,..,T_r; the first ellipsis has two dots.
is again regular.

In theorem 7.1. the normality assumptions as well as reducedness assumptions are quite superfluous, only lemme 7.2. is used (which should be stated in paragraph 1, maybe in a single proposition with the analoguous statement concerning normalisty, and properties $(R_k)$ and $(S_k)$ ). On the other hand the proof on page 32 seems to me \deleted{\unsure{u|i}ncom} a little short at one point; I tried to indicate some details in the margin.

As already stated in my notes, the statement of 9.4. is not the most natural one; one should use here the notion of tame ramification in the broad sense (just one divisor $D$ being given, whith normal crossings relatively to $S$). \deleted{This \unsure{w}} The proof would, I guess, require to use the results of sections 7 and 8 in a slightly different form (using the suitable variant of the notion of tame ramification, cf \deleted{\unsure{e|r}xample} particula\unsure{r}
% Only the left part of the final r survives at the right paper/image edge.
case b) in my notes). This however does not make any difficulty, in fact the proofs given in 7 and 8 work for any kind of ramification data, provided we assume a statement of the type 7.2. in paragraph 7 .

In 10.2. and 10.3. the significance of the result comes from the fact that the groups considred are just ${}_1(U,\ )$ resp. ${}_1(U,\ )^{(p)}$, and the theorems should be stated accordingly.
% Both group expressions have a blank base before subscript 1 and a blank second argument.

I hope you will be able to get along with these rather sketchy indications. Please \unsure{d\deleted{i}o|d\deleted{o}i} not hesitate to call back on me for whatever does not seem clear to you !

With best wishes,\\
% archive sheet number: 11
